\documentclass[11pt,a4paper]{article}
\usepackage[utf8]{inputenc}
\usepackage[T1]{fontenc}
\usepackage[spanish,es-noquoting]{babel}
\usepackage[margin=2.3cm]{geometry}
\usepackage{booktabs}
\usepackage{parskip}
\usepackage{amsmath}
\usepackage{caption}
\renewcommand{\arraystretch}{1.15}
\captionsetup{font=small,labelfont=bf}
\setlength{\parskip}{0.55em}

\begin{document}

\begin{center}
{\large\bfseries Informe de avance --- semana del 12 al 18 de agosto de 2026}\\[6pt]
{\small Para: Prof.\ Cristian Cappo \quad\textbullet\quad
De: Romina Alfonzo y Carlos Urdapilleta}\\[4pt]
{\footnotesize\itshape Detección de familias de ransomware en base a archivos encriptados
y notas de rescate}
\end{center}

\vspace{0.3em}
\hrule
\vspace{0.9em}

Esta semana se atendieron los pedidos de la reunión del 12/08 sobre el frente de archivos
cifrados. Todas las cifras provienen de corridas en el clúster del NIDTEC y quedaron
registradas con su manifiesto de parámetros y semillas.

\section*{1. Pedidos resueltos}

\textbf{El gráfico de bytes: la curva sí satura.} Se extendió la ablación de ventana de 512 a
1024, 2048 y 4096 bytes por extremo. El \textbf{macro-F1} por ventana es:

\begin{center}
\small
\begin{tabular}{lccccccc}
\toprule
Bytes leídos (cabecera + cola) & 128 & 256 & 512 & \textbf{1.024} & 2.048 & 4.096 & 8.192 \\
\midrule
macro-F1 & 0,790 & 0,792 & 0,851 & \textbf{0,905} & 0,904 & 0,903 & 0,901 \\
\bottomrule
\end{tabular}
\end{center}

El máximo está en 512+512 y a partir de ahí desciende levemente. Leer más de 1.024 bytes en
total no mejora el resultado: la información se agota antes.

Se agregó un control que no estaba previsto y conviene declarar. Al agrandar la ventana, los
archivos más cortos que ella se rellenan con ceros, y ese relleno codifica el tamaño del
archivo, que es una pista ajena al contenido. Por eso la curva se midió dos veces: sobre el
corpus completo (15.000 archivos) y sobre el subconjunto sin relleno ni solapamiento
(14.783 archivos, de al menos 8.192 bytes). Las dos curvas van paralelas, con el subconjunto
limpio unos 0,005 de macro-F1 por encima, de modo que \textbf{la forma de la curva no es un
artefacto del relleno}.

\textbf{Los bytes del medio no aportan.} Un bloque de 1.024 bytes tomado del centro del
archivo alcanza un \textbf{macro-F1 de 0,058}, con azar en 0,033. Sumado a los extremos, el
macro-F1 pasa de 0,905 a 0,903: no aporta y estorba levemente. La asimetría entre extremos es
marcada ---\textbf{cola sola 0,746} frente a \textbf{cabecera sola 0,348} de macro-F1---,
coherente con que once de las quince firmas binarias halladas sean sufijos y solo cuatro
prefijos. La marca la escribe el ransomware al final del archivo.

\textbf{Error incluido en los dos frentes.} Se repitieron las corridas canónicas sobre diez
semillas de muestreo, con hiperparámetros fijos, de modo que lo que se mide es la dispersión de
la estimación y no una nueva selección de modelo.

\begin{center}
\small
\begin{tabular}{ll}
\toprule
\textbf{Experimento} & \textbf{Resultado} \\
\midrule
2c --- aprendizaje sobre bytes & exactitud \textbf{0,912 $\pm$ 0,002} · macro-F1 \textbf{0,911 $\pm$ 0,001} \\
2b --- detección estructural, modo combinado & exactitud \textbf{0,932 $\pm$ 0,001} \\
3 --- notas, P1 (plantilla conocida) & macro-F1 \textbf{0,760 $\pm$ 0,029} \\
3 --- notas, P2 (variante nunca vista) & macro-F1 \textbf{0,435 $\pm$ 0,057} \\
\bottomrule
\end{tabular}
\end{center}

\textbf{BLACKBASTA incorporada.} Los dos frentes se miden ahora sobre \textbf{30 familias},
con azar en 0,033. Desaparece la asimetría 30/29 que había que aclarar en cada tabla.

\textbf{Validación separada en train y test.} Ya estaba implementada y esta semana quedó
verificada en el código: el clasificador de bytes usa validación cruzada \textbf{anidada}
---la búsqueda de hiperparámetros se ajusta dentro del pliegue de entrenamiento y se evalúa en
el pliegue externo, dando 0,891 de exactitud--- y el de notas usa \texttt{StratifiedGroupKFold}
agrupando por plantilla, de modo que bajo el protocolo P2 ninguna variante del conjunto de
prueba aparece en el de entrenamiento.

\section*{2. Hallazgos nuevos}

\textbf{Se identificó la marca de BADRABBIT, que figuraba como familia sin marca.} Escribe la
palabra \texttt{encrypted} codificada en UTF-16 al final de los archivos que cifra, en
\textbf{965 de 1.000 archivos (96,5\,\%)}. El detector no la encontraba porque exigía que la
marca apareciera en \emph{todos} los archivos de la muestra, y con un 3,5\,\% de excepciones la
probabilidad de que una muestra de 50 archivos no contuviera ninguna era de apenas 0,17.
\textbf{Las familias sin marca detectable bajan de cuatro a dos: NOTPETYA y SUNCRYPT.}

\textbf{Se corrigió un defecto metodológico del detector.} El criterio de unanimidad byte a
byte se reemplazó por un umbral de mayoría declarado en 0,90, y cada marca se reporta con su
cobertura. Medido sobre diez semillas, el criterio nuevo \textbf{reduce el desvío catorce
veces} (exactitud del modo combinado 0,932 $\pm$ 0,001 frente a 0,900 $\pm$ 0,016) y además
\textbf{sube la media 0,032}. Con el criterio anterior, el número de familias con marca
oscilaba entre 26 y 28 según el sorteo; con el nuevo son 28 de 30 en las diez semillas.

\textbf{Queda explicado por qué fallan seis familias y no dos.} Cruzando el tipo de marca que
deja cada familia con su macro-F1 en el clasificador:

\begin{center}
\small
\begin{tabular}{lcc}
\toprule
\textbf{Marca que deja la familia} & \textbf{Familias} & \textbf{Con macro-F1 $\geq$ 0,98} \\
\midrule
Firma binaria y extensión & 15 & 15 \\
Solo firma binaria & 2 & 1 \; (BADRABBIT queda en 0,978) \\
\textbf{Solo extensión} & 11 & \textbf{7} \\
Sin marca alguna & 2 & 0 \\
\bottomrule
\end{tabular}
\end{center}

Las cuatro que fallan del grupo «solo extensión» son JIGSAW, DARKSIDE, CRYPTOLOCKER y
WASTEDLOCKER. Como el clasificador de bytes \textbf{no usa el nombre ni la extensión}, se queda
sin la única marca que esas familias dejan. Sumadas a NOTPETYA y SUNCRYPT dan las seis familias
difíciles, que son las mismas de la corrida anterior, ahora con desvío: NOTPETYA
0,394 $\pm$ 0,023 · JIGSAW 0,439 $\pm$ 0,014 · DARKSIDE 0,603 $\pm$ 0,011 · CRYPTOLOCKER
0,605 $\pm$ 0,016 · WASTEDLOCKER 0,628 $\pm$ 0,013 · SUNCRYPT 0,745 $\pm$ 0,016.

\textbf{La dispersión de cada frente es, en sí misma, un resultado.} El frente de archivos
tiene un desvío de $\pm$ 0,001 de macro-F1 sobre 15.000 muestras; el de notas, $\pm$ 0,057
sobre 95 plantillas. Es una diferencia de un factor cercano a cuarenta, y refuerza por una vía
independiente lo que ya indicaban los tres resultados negativos de optimización: \textbf{el
techo del frente de notas lo impone la cantidad de datos, no el método.}

\textbf{Validación externa de las firmas.} Se recuperaron las anotaciones de la evaluación
manual de ID Ransomware y seis familias presentan coincidencia exacta entre la firma que el
detector descubre por su cuenta y los bytes que la herramienta reporta: WANNACRY, LORENZ, MAZE,
GANDCRAB, MEDUZALOCKER y TESLACRYPT. Son dos derivaciones independientes que llegan a los
mismos bytes.

\textbf{Verificación de integridad del corpus.} Se detectaron 12 archivos sin cifrar ---JPEG en
claro, confirmado por sus bytes de cabecera--- dentro de la carpeta de CERBER. Excluirlos mueve
la exactitud del Experimento 2c en 0,003, dentro de la dispersión entre semillas. No se
determinó si se trata de un artefacto del empaquetado del conjunto o de archivos que la familia
no cifró.

\section*{3. Observación sobre el conjunto de datos mayor}

Se plantea una salvedad antes de invertir en la descarga. El artículo de NapierOne indica que
cada familia se obtuvo \textbf{ejecutando una sola vez} una muestra del ransomware en una
máquina preparada. Una escala mayor del conjunto aporta más archivos \emph{de la misma
campaña}, de modo que sirve para estabilizar las métricas pero no para evaluar generalización
entre campañas. La limitación de una campaña por familia no se resuelve por volumen.

\section*{4. Próximos pasos}

\textbf{Frente de notas, en curso.} Curva de aprendizaje del rendimiento contra cantidad de
plantillas por familia, que responde con un número cuántas notas harían falta, y análisis de los
marcadores compartidos entre plantillas de una misma familia para intentar detectar una variante
no vista.

\textbf{Cruce con el nombre del archivo (elemento de acción 2).} Diseñado, sin correr. Se
propone reportar tres condiciones comparables: solo bytes (0,911 de macro-F1, la referencia),
bytes más la \emph{forma} del nombre ---largo de la extensión, si es hexadecimal o pronunciable,
si conserva el nombre base, si incorpora un correo o un identificador--- y bytes más la
extensión literal. La segunda es la que puede sostenerse entre campañas; la tercera se
reportaría como cota superior declarada, porque la extensión identifica la campaña y no la
familia.

\textbf{Consulta pendiente.} El \emph{majority voting} combinando notas y archivos tiene un
obstáculo que conviene discutir: no existen muestras pareadas. Las notas provienen de
repositorios públicos y los archivos cifrados de NapierOne, de modo que no hay un incidente del
que se tengan los dos artefactos. El esquema puede proponerse como arquitectura de despliegue,
pero no evaluarse con los datos disponibles.

\vspace{1em}
\hrule
\vspace{0.5em}
{\footnotesize Los experimentos de mayor costo computacional se ejecutaron en el clúster del
Núcleo de Investigación y Desarrollo Tecnológico (NIDTEC) de la Facultad Politécnica, UNA,
financiado por el proyecto LABO16-167 del programa PROCIENCIA (CONACYT).}

\end{document}
